\section{\filetotitle}
\textit{Commentary on: \fullcite{10.1021/acs.jctc.5c01176}}

The second published work of this thesis focuses on curvature-driven surface hopping algorithms from article \textit{Curvature-Driven Nonadiabatic Dynamics without Explicit Nonadiabatic Couplings} in \textit{Journal of Chemical Theory and Computation}.

When performing nonadiabatic dynamics simulations, one has several options how to treat the nonadiabatic transitions between electronic states. The most common approach is to use surface hopping methods, where classical nuclear trajectories evolve on potential energy surfaces and can hop between them based on certain criteria. The most widely used surface hopping method is the \acrfull{fssh} method, described in Sec~\ref{sec:fewest_switches}, which relies on the calculation of \acrfull{tdc} to propagate propagate the electronic wavefunction in adiabatic basis. Details about the \acrshort{tdc} can be found in Sec~\ref{sec:time_derivative_coupling} and the adiabatic basis is described in Sec~\ref{sec:adiabatic_vs_diabatic_representation}. From quantum-chemical calculations, the \acrshort{tdc} is usually obtained using \acrfull{nacv} (Eq~\eqref{eq:tdc_nacv}) that are computationally expensive to calculate and not available for all electronic structure methods. The question arises whether it is possible to perform nonadiabatic dynamics simulations without the need to calculate \acrshort{nacv} at all.

One of the fairly known approaches to avoid the calculation of \acrshort{tdc} is to use the \acrfull{lzsh}, described in depth in Sec~\ref{sec:landau_zener}, which estimates the hopping probabilities based on the exact solution of a two-level system with linearly varying diabatic energies and constant diabatic coupling. The \acrshort{lzsh} method has been successfully applied in various nonadiabatic dynamics simulations, but it is not as popular as the \acrshort{fssh} method due to its theoretical limitations and fairly recent applications to simluations in adiabatic basis. Since most of the quantum-dynamical codes already provide \acrshort{fssh} implementation, it would be beneficial to have a method that can be easily incorporated into existing \acrshort{fssh} frameworks without the need to calculate \acrshort{nacv} from \textit{ab initio} methods.

The study presented in this article aims to evaluate the performance of the newly proposed curvature-driven $\kappa$FSSH method in comparison to the traditional \acrshort{fssh} and \acrshort{lzsh} methods. The theory behind the $\kappa$FSSH method is described in detail in Sec~\ref{sec:baeck_an}. The surface hopping methods are tested on several one-dimensional model systems, including single avoided crossing, three-state crossing and models exhibiting trivial crossings. After the on-dimensional models, the study applies the algorithm to more realistic systems as 8-dimensional model of uracil cation or nonadiabatic dynamics of stilbene, azobenzene or cyclobutanone. In the simulations, the quantities of interest are the state populations and the magnitude if the \acrshort{tdc}. The populations obtained from the surface hopping methods are compared to the exact quantum dynamics results (Sec~\ref{sec:split_operator_method}), while the \acrshort{tdc} values are compared to those calculated from the \acrfull{npi} algorithm described in Sec~\ref{sec:norm_preserving_interpolation}. The study also introduces a modified version of the $\kappa$FSSH method, called $\lambda$FSSH, which should theoretically be more accurate in predicting the \acrshort{tdc} values (Eq~\eqref{eq:lambda_tdc}).

Since the algorithms are tested on potentials where a three-state crossing occurs, the study also develops an algorithms to use the \acrshort{lzsh} method in such cases, which is not straightforward due to the fact that the \acrshort{lzsh} method is originally designed for two-level systems. The suggested algorithms for applying the \acrshort{lzsh} method to three-state crossings are described in detail at the end of Sec~\ref{sec:landau_zener}.

The results from the simulations on a single avoided crossing model show that the $\kappa$FSSH already performs significantly worse than the traditional \acrshort{fssh} and \acrshort{lzsh} methods in terms of state populations. Inspecting the \acrshort{tdc} in this case reveals that the $\kappa$FSSH method produces a narrower peak compared to the reference \acrshort{npi} values, which likely contributes to the poorer performance in state populations. In the three-state crossing model, the performance of the $\kappa$FSSH method is even more unsatisfactory, with the state populations deviating significantly from the exact results. The \acrshort{tdc} values also show considerable discrepancies compared to the reference values, indicating that the curvature-driven approach may not be suitable for accurately capturing nonadiabatic transitions in more complex systems. Moreover, $\kappa$FSSH predicts non-zero \acrshort{tdc} values even between states, where no coupling should exist, leading to unphysical transitions.

Moving on to the trivial crossing models, the $\kappa$FSSH method again fails to reproduce the correct state populations, while the traditional \acrshort{fssh} and \acrshort{lzsh} methods perform without issues. The \acrshort{tdc} values from the $\kappa$FSSH method exhibit significant deviations from the reference values, further highlighting the limitations of the curvature-driven approach in capturing nonadiabatic dynamics accurately. Investigating the \acrshort{tdc} values in detail reveals that the $\kappa$FSSH method produces significantly smaller peaks compared to the reference values, which explains the poor performance in state populations, since the \acrshort{tdc} should be infinite at the trivial crossing point.

Applying the algorithm to the vibronic coupling model of uracil cation, the error of the $\kappa$FSSH method in state populations is less pronounced. \acrshort{fssh}, \acrshort{lzsh} and $\kappa$FSSH all perform reasonably well in this case. This suggests that the curvature-driven approach may be more suitable for systems with higher dimensionality, since the additional degrees of freedom can help to average out the inaccuracies in the \acrshort{tdc} values. However, aplying the $\kappa$FSSH method to the real systems of stilbene, azobenzene and cyclobutanone reveals that the method still significantly underperforms compared to the traditional \acrshort{fssh} and \acrshort{lzsh} methods. The state populations obtained from the $\kappa$FSSH method deviate considerably from the reference results, but the reason for this behavior is different than in the previous cases. In these realistic systems, the electronic structure calculations often yield discontinuous potential energy surfaces due to rotations of orbitals or changes in the electronic wavefunction character. These discontinuities lead to erratic \acrshort{tdc} values for the $\kappa$FSSH method, which in turn results in spurious hops and inaccurate state populations.

In summary, the study presented in this article demonstrates that while the curvature-driven $\kappa$FSSH method offers a promising approach to avoid the calculation of \acrshort{nacv}, it falls short in accurately capturing nonadiabatic dynamics across a range of model systems and realistic molecules. The traditional \acrshort{fssh} and \acrshort{lzsh} methods consistently outperform the $\kappa$FSSH method in terms of state populations and \acrshort{tdc} values. The limitations of the curvature-driven approach are particularly evident in systems with trivial crossings and discontinuous potential energy surfaces. Future research may focus on refining the curvature-driven algorithms to improve their accuracy and reliability in nonadiabatic dynamics simulations. On the other hand, the $\kappa$FSSH method could still find its use in specific applications where the calculation of \acrshort{nacv} is prohibitively expensive and the \acrshort{pes} are smooth enough to avoid the issues observed in this study.
